\color{unimain}{\subsection{Защита от потери охлаждения}}
\color{black}

\begin{enumerate}

\item
Защита от потери охлаждения направлена на предотвращение перегрева и повреждения шунтирующего реактора в результате частичного или полного отказа системы охлаждения.

\item
Функция защиты при потере охлаждения контролирует уровень тока ВВ, токовый орган выводится программной накладкой \textbf{«РТ ВВ ЗПО»}. При необходимости можно выводить отдельные ступени программными накладками \textbf{«Работа ЗПО-1»}. Программная накладка \textbf{«Работа ЗПО»} выводит выходные сигналы. Присутствует возможность выводить отключающую цепочку программной накладкой \textbf{«Откл ЗПО»}.

\item
Защита содержит три ступени, каждая из которых может быть выполнена с контролем температуры масла, которая вводится программной накладкой и уровня тока. Ввод контроля перегрева масла, а также выбор соответствующего уровня тока осуществляется программными накладками \textbf{«ЗПО-1 конт.ДТм»} и \textbf{«ЗПО-1 конт.РТ»} соответственно. Программная накладка \textbf{«ЗПО-1 конт.РТ»} имеет три положения:
\begin{itemize}
\item \textit{«Откл»} -- выводит ступень; 
\item \textit{«От РТ1»} -- от первого токового реле;
\item \textit{«От РТ2»} -- от второго токовго реле. 
\end{itemize}

Каждая из ступеней имеет индивидуальную регулируемую задержку на срабатывание \textbf{«Тср ЗПО-1»}.

\item 
Уставки ЗПО приведены в таблице \ref{app_set:zpo}. Алгоритм работы ЗПО приведен на рисунке \ref{pict_zpo:ZPO}.

\medskip
\begin{figure}[H]
\centering
\includegraphics[width=0.7\textwidth,height=0.7\textheight,keepaspectratio]{\fbpath/11.01 ШР/171.1151 ЗПО/img/zpo.pdf}
\captionof{figure}{Функциональная схема логики ЗПО}\label{pict_zpo:ZPO}
\end{figure}


\end{enumerate}